% a06_BALtorus_d6: part of Step 3 of the proof of thm:a06:T1 (input from sec/a06_BALtorus_c.tex).
% Group (D): Ba3 Prop. 1 at L = 5, first half: localisation conventions and the regrouping after separated lines are cut (p. 432).

\subsubsection*{Step 3(D$_3$): \cite[Prop.~1]{Ba3} at $L=5$, the regrouping after the cut}

The proof of \cite[Prop.~1, (1.33)]{Ba3} on \cite[pp.~420--445]{Ba3} consists of twenty sketched steps and seven further sketched places. Most of them are direct computations on the printed definitions (the slices \cite[(2.6), (2.10)--(2.12)]{Ba3}, the summation by parts \cite[(2.8)]{Ba3}, the induction \cite[(2.16)]{Ba3}, the generalised graphs of \cite[Prop.~2.2]{Ba3}, the estimate \cite[(2.17)]{Ba3}, the Ward--Takahashi identities of \cite[pp.~430--431]{Ba3} for full lines, the estimate \cite[(3.1)]{Ba3}, the cube steps of \cite[p.~433]{Ba3}, the norms \cite[(1.32)]{Ba3} at $B_0$ and at $\tilde B$), whose written form is a transcription with the $L=5$ constants of \cref{thm:a06:BALA}. Two steps need more than a transcription: the regrouping into renormalised classes after the separated lines are cut (\cite[p.~432]{Ba3}: ``if we remove a line in any graph of the class, then for the obtained graph we can find in this class all graphs necessary to form a new renormalized class''), and the renormalisation of the divergent classes (\cite[pp.~434--445]{Ba3}). They are proved here and in Step~3(D$_4$); the conventions are Ba{\l}aban's, made explicit.

\paragraph{Conventions.} (CV) Degrees are computed in Ba{\l}aban's count: each of \cite[(1.13), (1.14), (1.15)]{Ba3} is a vertex with its own degree (\cite[p.~423 (v), p.~414]{Ba3}); the two factors of one product \cite[(1.18)]{Ba3} share the block point $y$ and are joined in the bound function by a contact, a top-scale line of dimension $0$ whose kernel is the indicator that both summation points lie in $B^k(y)$. (Loc) Label boxes are closed grid boxes of side $\le2$ on the block-corner grid (\cite[p.~420]{Ba3}: ``for each leg of the vector field a smooth partition of unity \dots\ a support of each function is contained in a cube with sides of length 2''); a label is attached to each leg of a vector field (the support box of its partition function $H_y\subset y+[-1,1]^d$), to each \cite[(1.6)]{Ba3}/\cite[(1.7)]{Ba3} vertex (its unit cube) and to each block factor ($B^k(y)$), and every vertex lies in the intersection of the boxes of its labels. The number of label boxes of a graph is $V_{\mathrm{lab}}\le(2\bar n+1)V_{\mathrm{Bal}}$, $V_{\mathrm{Bal}}\le2\bar n(4N_A+5)$, constants of $\bar n$ only; put $l_N\defeq2V_{\mathrm{lab}}(\bar n)$. (Agr) The counterterm graph $T_\gamma G$ of a two-scalar-leg subgraph $\gamma$ is $G$ with $\gamma$'s two external legs attached at one point $x$ \cite[p.~417]{Ba3}, $x$ the leg point of the anchor vertex $v_a$ of $\gamma$; $\gamma$'s internal vertices keep their labels (\cite[p.~420]{Ba3}, the localisation is performed ``for the terms of $\delta m^2_k(\dots)$, $E_k(\dots)$ also''), so $G$ and $T_\gamma G$ have the same label multiset (\cref{conv:a06d:LOC}(c)). (R-W$'$) Two graphs are ``in agreement'' \cite[p.~432]{Ba3} if they carry the same label multiset; a \emph{per-label class} of a given type and order is the set of all graphs of that type and order with a given label multiset, equivalently (with the partition functions written as multipliers $t_yH_y$ on the vector legs and $t_Q\mathbf 1_Q$ on the \cite[(1.6)]{Ba3}/\cite[(1.7)]{Ba3} vertices) the coefficient of one monomial $\prod t$ in the expansion of \cite[(1.5)]{Ba3}; this is how \cite[p.~444]{Ba3} evaluates ``a product of the values of all the localization functions at the point $x'$'' and then takes the graph ``with the summations unrestricted''. (H$_A$) For an external field with a leg that can sit at several vertices of a class (the counterterm merge moves it to the anchor, leg exchanges move it too), the localisation $h$ of \cite[p.~420]{Ba3} (``$h=1$ on box$(v)$ and $h=0$ outside some neighborhood of box$(v)$'') is taken $=1$ on the label boxes of every vertex at which the leg sits in some term of the class and $0$ outside their $\frac14$-neighbourhood; for moves inside a cluster this is a neighbourhood of size $\le l_N$, and for moves across clusters the extra box is paid by \cref{lem:a06d:S25}(a). It costs nothing to the consumers, which sum over localisations \cite[p.~421]{Ba3} or use fields localised at one point.

\begin{definition}[clusters; the cluster cut rule]\label{def:a06d:CCR}
Let $\tau$ be a term of a per-label class and $\mathcal B$ its multiset of label boxes, common to the class. The \emph{clusters} are the connected components of the graph on $\mathcal B$ in which two boxes are joined when they intersect; $U_K$ is the union of the boxes of the cluster $K$. Each end of a line of $\tau$ lies at a vertex, all of whose labels contain it, so they intersect and lie in one cluster $K(\mathrm{end})$ (for a counterterm vertex, the cluster of $v_a$). The \emph{cut set} $S(\tau)$ is the set of lines whose two ends lie in different clusters. A cut line becomes the external field $h_KGh_{K'}$ (or its vector analogue), with $h_K$ smooth, $h_K=1$ on the $\eta$-neighbourhood of $U_K$, $\supp h_K$ in the $\frac14$-neighbourhood of $U_K$ and $|\partial^\eta h_K|,\mathrm{Lip}(\partial^\eta h_K)\le c(d)$ (e.g.\ $h_K=1-\prod_Q(1-\varphi_Q)$ over the unit grid cubes $Q$ of $U_K$, $\varphi_Q$ a fixed bump). This is Ba{\l}aban's threshold $1$ \cite[p.~432]{Ba3}, applied to clusters instead of single lines.
\end{definition}

\begin{lemma}[properties of the cut rule]\label{lem:a06d:S22}
\textup{(a)} If $K\ne K'$ then $\dist_\infty(U_K,U_{K'})\ge1$, and points of $\supp h_K$ and $\supp h_{K'}$ are at distance $\ge\frac12$. \textup{(b)} $\operatorname{diam}_\infty U_K\le2|K|\le l_N$. \textup{(c)} Every kept line joins two vertices of one cluster, so every kept-connected component lies in one $U_K$. \textup{(d)} Let $\gamma$ be a subgraph all of whose label boxes lie in one cluster $K$; any operation that changes only the vertices of $\gamma$ at which lines are attached (the counterterm merge of \cite[p.~417]{Ba3}, the exchange of two external legs of $\gamma$), or replaces $\gamma$ by another graph with the same labels, leaves $S$ unchanged. \textup{(e)} A connected subgraph whose vertices lie in at least two clusters contains a cut line and is disconnected by $S$, each piece lying in one cluster. \textup{(f)} A subgraph whose vertices lie in one cluster has no cut line.
\end{lemma}

\begin{proof}
(a) Two grid boxes either intersect or are at sup-distance $\ge1$; boxes of different clusters do not intersect; the collars add $\le\frac14$ on each side. (b) Two boxes of $K$ are joined by a chain of at most $|K|-1$ intersecting boxes of side $\le2$. (c) is the definition. (d) The moved ends stay at vertices of $\gamma$, whose labels are in $K$, and the other ends do not move; a replacement with the same label multiset does not change $\mathcal B$, hence not the clusters. (e) Kept lines are intra-cluster, so a kept path cannot leave a cluster. (f) is immediate.
\end{proof}

Let $G$ be a labelled graph and $\mathcal C=R(G)$ its class: the forest sum over $\mathcal D$, the partnered subgraphs of $G$ (the one-particle-irreducible two-scalar-leg subgraphs with mass counterterms, \cite[p.~430]{Ba3}, and the vacuum components with $E_1$, \cite[p.~432]{Ba3}), closed under leg exchange of $2\phi$ and $\phi\phi A$ subgraphs and under per-label completion of $2A$ and $\phi\phi A$ subgraphs. Split $\mathcal D=\mathcal D_0+\mathcal D_1$, $\mathcal D_0\defeq\{\gamma:\text{all label boxes of }\gamma\text{ in one cluster}\}$. For a forest $F=F_1+F_0$ ($F_1\subset\mathcal D_1$, $F_0\subset\mathcal D_0(F_1)$, the elements of $\mathcal D_0$ nested in or disjoint from every element of $F_1$), the term $\tau_F\defeq\prod_{\gamma\in F}(-T_\gamma)G$ is cut by its own cut rule.

\begin{theorem}[the regrouping after the cut]\label{thm:a06d:S23}
\textup{(i)} For fixed $F_1$ the cut set $S(\tau_{F_1+F_0})$ is the same for every $F_0$; call it $S(F_1)$. \textup{(ii)} $\sum_FE(\tau_F)=\sum_{F_1}R_{\mathcal D_0(F_1)}[E(\mathrm{Cut}_{S(F_1)}Y(F_1))]$, $Y(F_1)\defeq\prod_{\gamma\in F_1}(-T_\gamma)G$, each $F_0$-operation acting at its place (inside a counterterm coefficient when nested in an element of $F_1$). \textup{(iii)} Every element of $\mathcal D_1$ is disconnected by $S(\emptyset)$. \textup{(iv)} Every connected divergent subgraph of $\mathrm{Cut}\,Y(F_1)$ (kept lines only) is an element of $\mathcal D_0(F_1)$, renormalised by $R_{\mathcal D_0(F_1)}$; or a $2A$ or $\phi\phi A$ subgraph inside one cluster whose per-label partners are terms of the same block (\cref{thm:a06d:S28}); or a \cite[(2.4)]{Ba3} subgraph (no partner; treated by the summation by parts \cite[(2.8)]{Ba3}); or a subgraph containing a counterterm vertex of $F_1$, which by \cref{lem:a06d:S25} is convergent unless it is a vacuum graph or has one external vector leg (Ba{\l}aban's exception list, \cite[p.~432]{Ba3}). So each block is a renormalised class of its cut graph (\RENorb, \cref{hyp:k01:RENorb}, is the same regrouping with the slice decomposition in place of the cut).
\end{theorem}

\begin{proof}
(i) The $F_0$-operations move ends inside clusters (\cref{lem:a06d:S22}(d)). An element $\gamma_0\in\mathcal D_0(F_1)$ contains no line whose end is moved by an element $\gamma\in F_1$: such a line is external to $\gamma$ and attached at a vertex of $\gamma$, so $\gamma_0$ would share a vertex with $\gamma$, hence be nested with it by compatibility; $\gamma_0$ inside $\gamma$ contains no external line of $\gamma$, and $\gamma_0$ containing $\gamma$ would contain $\gamma$'s cut line, against $\gamma_0\in\mathcal D_0$. For vacuum elements of $F_1$ nothing moves: the vacuum graph and its counterterm graph (same graph, same labels, \cref{conv:a06d:LOC}(b)) are cut identically. (ii) Every forest of $\mathcal D$ splits uniquely as $F_1+F_0$ with $F_1$ a forest of $\mathcal D_1$ and $F_0$ a forest of $\mathcal D_0(F_1)$, and conversely; by (i) the cut is the same for all $F_0$ at fixed $F_1$, and $\mathrm{Cut}$ commutes with the $F_0$-operations because they act inside clusters and the elements of $\mathcal D_0$ have no cut line (\cref{lem:a06d:S22}(f)). (iii) \cref{lem:a06d:S22}(e). (iv) A divergent connected kept subgraph $H$ of $\mathrm{Cut}\,Y(F_1)$ without a counterterm vertex of $F_1$ lies in one cluster (\cref{lem:a06d:S22}(c)) and avoids the internal vertices of the elements of $F_1$ (these sit inside the counterterm coefficients). If $H$ is $2\phi$ it belongs to $\mathcal D$, lies in one cluster and is disjoint from $F_1$, so $H\in\mathcal D_0(F_1)$; if $H$ is $2A$ or $\phi\phi A$, \cref{thm:a06d:S28}; \cite[(2.4)]{Ba3} needs no partner; subgraphs through a counterterm vertex of $F_1$ are \cref{lem:a06d:S25}.
\end{proof}

\begin{corollary}[the cut kernels]\label{cor:a06d:S24}
In every term of every block, every cut line joins two different clusters, so its kernel is evaluated on $\supp h_K\times\supp h_{K'}$, at distance $\ge\frac12$. The estimate \cite[(3.1)]{Ba3} holds with $h_K,h_{K'}$ in place of $h,h'$: $\norm{h_KG_k(\Omega,\tilde B)h_{K'}}_{1,\alpha}\le Ce^{-(\delta_1/2)(\dist(U_K,U_{K'})-1/2)}$, its proof using only the slice bounds at the points of the supports, the separation, and $|\partial h|,\mathrm{Lip}(\partial h)\le c_h$; the points used lie within $\frac14+\eta$ of vertex label boxes, at the depth the hypotheses supply ($R\ge R_0+5$). The tree factor of \cite[(1.33)]{Ba3} satisfies $d(\{\square(v)\})\le\sum_{\text{spanning tree of cut kernels}}\dist(U_K,U_{K'})+\sum_K\operatorname{diam}U_K$, and the second sum is $\le l_N^2$, a constant.
\end{corollary}

With the printed per-line threshold $\dist(\square(v),\square(v'))\ge1$ and the counterterm leg move of \cite[p.~417]{Ba3}, a line from the moved leg of $\gamma$ to a vertex $w$ could be cut in $G$ and short in $T_\gamma G$, or conversely, and then either the cut kernel is not $O(1)$ or the cut does not commute with $T_\gamma$; the cluster rule removes this (the example of three boxes $[0,1],[1,2],[2,3]$ on one axis is one cluster, and the line is kept in $G$ and in $T_\gamma G$), with Ba{\l}aban's threshold $1$ and no larger cut threshold.

\begin{lemma}[disconnected counterterms are finite vertices]\label{lem:a06d:S25}
Let $\gamma\in\mathcal D_1$ (so disconnected by $S$), with anchor $v_a$, and let $c(x)$ be the coefficient of $T_\gamma$ in $Y(F_1)$ with $\gamma$'s labels fixed: $c(x)=\sum_{x'}$ of the amputated kernel of $\gamma$, renormalised by the inner forests, with its lines of $S(F_1)$ replaced by cut kernels. Assume \cite[(1.33)]{Ba3} for all classes with fewer vertices than $G$ (the induction on the number of vertices that the classes already use, \cite[p.~431]{Ba3}: ``our definition of the classes $G_{\mathrm{ren}}$ is in fact inductive''). Then \textup{(a)} $|c(x)|\le Ce_k^{d_v(\gamma)}\lambda_k^{d_s(\gamma)}\exp(-\delta'd(\text{labels of }\gamma))$ uniformly in $x,k,\eps,\eta,\Omega,\Omega_1,\Omega_2,\tilde B$ and the torus; \textup{(b)} $T_\gamma$ is therefore a \cite[(1.7)]{Ba3} vertex with a finite counterterm, of degree $D(v)=2$ (\cite[p.~423 (ii)]{Ba3}); every connected subgraph containing it has $D\ge1$ if $E\ge2$ and $D\ge0$ if $E=1$, with equality only for the tadpole graph, and vacuum graphs are treated by the vacuum rule of \cite[p.~432]{Ba3}; \textup{(c)} $c$ is an even polynomial in $q$, so the tadpole with this vertex vanishes after the gauge-out of $B_0$ \cite[p.~433]{Ba3}.
\end{lemma}

\begin{proof}
(a) The pieces of $\gamma|_{S(F_1)}$ lie in single clusters, each with fewer vertices than $G$ since $\gamma$ has at least two pieces; each piece with its inner renormalisation is a class of lower order whose external fields are the cut kernels, two-variable external fields in the sense of \cite[p.~419]{Ba3}; \cite[(1.33)]{Ba3} for each piece, the product over pieces, and \cref{cor:a06d:S24} for the cut kernels give (a), the decay coming from the cut kernels between the clusters of $\gamma$. (b) A \cite[(1.7)]{Ba3} vertex with a bounded coefficient has $D(v)=d-(d-2)=2$, and the degree count of \cite[(2.1)--(2.3)]{Ba3} applies with $\kappa=2$ at that vertex. (c) $\gamma$ has no external vector leg; every closed scalar index loop of $\gamma$ carries an even number of vector legs or vanishes ($\Tr q^{\mathrm{odd}}=0$), and the scalar path between $\gamma$'s two legs then carries an even number as well, since the total number of vector-leg ends of $\gamma$ is twice its number of photon lines.
\end{proof}

\begin{lemma}[per-label gauge invariance]\label{lem:a06d:S26}
Let $Y$ be a periodic box or $\xi\Z^d$, $X$ any bond field built from $A'$ and external fields (e.g.\ $X=\sum_yt_yH_yegA'+\tilde A$-type legs with their multipliers), $F$ any polynomial in the $|\phi(x)|^2$ with $x$-dependent coefficients (e.g.\ $\sum_Qt_Q\lambda\mathbf 1_Q|\phi|^4$, the \cite[(1.7)]{Ba3} vertices with any coefficients, the restriction to $\Omega_1$), and $Z(A;t)\defeq\int d\mu(A')\int d\phi\exp[-\frac12\langle\phi,(-\Delta_{A+X}+m^2)\phi\rangle-F(\phi)]$ with any Gaussian measure $d\mu(A')$. Then $Z(A+d\lambda;t)=Z(A;t)$ for every periodic (on $\xi\Z^d$: finitely supported) $\lambda$ and every $t$. Consequently every $t$-monomial coefficient of every $A$-Taylor coefficient of $\log Z$ satisfies the Ward--Takahashi identities of \cite[pp.~430--431]{Ba3}; in particular the two-vector-leg kernel $\Pi^{\mathrm{lab}}$ of a per-label class obeys $\sum_{y'}\Pi^{\mathrm{lab}}_{\mu\nu}(y,y')=0$.
\end{lemma}

\begin{proof}
$\phi=V_\lambda\psi$: the covariant Laplacian is covariant for the total field, $-\Delta_{A+d\lambda+X}=V_\lambda(-\Delta_{A+X})V_\lambda^*$ for every $X$, because only the total field enters $U$; $F$ is invariant, being a function of the $|\phi(x)|^2$; the $A'$-integral is untouched since the transformation acts on $A$ only. The identities follow coefficient by coefficient in $t$, and by (R-W$'$) the per-label class is one $t$-monomial.
\end{proof}

The identity is for complete per-label classes with full propagator lines, before any scale decomposition; it is never applied to a single ordering or scale assignment. In Step~3(D$_4$) it is applied only to the all-$C^\xi$ part of a regrouped local coefficient (full lines $C^\xi$ on $\xi\Z^3$), the remainders being bounded explicitly.

\begin{lemma}[{multipliers to the anchor; \cite[p.~444, (3.36)]{Ba3}}]\label{lem:a06d:S27}
For a node with anchor $\hat x$ and a smooth multiplier $H$ (partition function, $g_k$) at a vertex $x_i$ of the node, $H(x_i)=H(\hat x)+[H(x_i)-H(\hat x)]$ with $|H(x_i)-H(\hat x)|\le c|x_i-\hat x|$; the second term carries a factor $|x_i-\hat x|$ on a path inside the node, so its degree gains $1$. After all smooth multipliers of the node (including those of its external legs) are moved to $\hat x$, the remaining sums over the node's internal positions are unrestricted, except for sharp indicators (\cref{lem:a06d:S46}).
\end{lemma}

\begin{theorem}[the Ward--Takahashi and $\phi\phi A$ classes at fixed localisation]\label{thm:a06d:S28}
Let $H$ be a $2A$ or $\phi\phi A$ subgraph of a term of a per-label class. \textup{(a)} If the label boxes of $H$ lie in one cluster, every per-label partner $\tilde H$ (same type, order, label multiset) lies entirely in that cluster and has no cut line, and so does every leg-exchange partner (\cref{lem:a06d:S22}(d), (f)); the per-label completion is therefore a sum of terms of one block, and at the renormalisation step the local coefficient of the completed node is the per-label class sum, to which \cref{lem:a06d:S26} applies ($2A$ zeroth moment, after \cref{lem:a06d:S27} has moved the multipliers of $H$'s external legs to the anchor), as do \cref{lem:a06d:S27,lem:a06d:S46} with the reflection identities of \cite[(3.30)]{Ba3} ($\phi\phi A$ zeroth moment, $2A$ first moment). \textup{(b)} If they span at least two clusters, every partner is disconnected by $S$ (\cref{lem:a06d:S22}(e)) and no identity is needed.
\end{theorem}

\begin{proof}
\cref{lem:a06d:S22}.
\end{proof}

At fixed vertex localisations of all vertices of $H$ the zeroth-moment sum would be restricted to boxes and would not vanish (boundary layers of width $L^{j_0}\eta$); Ba{\l}aban localises legs, not vertices \cite[p.~420]{Ba3}, and moves the localisation functions to one point before the identity is used \cite[p.~444]{Ba3}, and with (R-W$'$) the identity holds label by label.

Consequences for the neighbouring steps: under the cluster rule the kept components lie in clusters (\cref{lem:a06d:S22}(c)), and the class-component cube of \cite[p.~432]{Ba3} has side $\le l_N$, common to the class; kept lines may join boxes at distance up to $l_N$, so \cite[(3.2)]{Ba3} for components carries the constant $e^{\delta l_N}$. The cut on the cube after the $\tilde B'$-expansion \cite[p.~433]{Ba3} is the cluster rule on the label boxes of the expanded class; every component with an old vertex lies in the cluster of an old vertex, of diameter $\le2V^\sharp_{\mathrm{lab}}$ ($V^\sharp_{\mathrm{lab}}$ the label count after the expansion), so the side of the localisation cube of a connected renormalised class after the expansion is $l^\sharp\defeq l_N+4V^\sharp_{\mathrm{lab}}(\bar n)$, which fixes $r^\flat$ in \cref{thm:a06:BALA}(a) through $r^\flat\ge3l^\sharp+R_0+12$. Chain components lie in one cluster, so their exponential factor is $\ge e^{-2\delta V^\sharp_{\mathrm{lab}}}$.
